% =====================================================================
%  TCC1 — PROJETO DE PESQUISA — arquivo de partida
%  Modelo ABNT do ICT/Unifei — segue a NBR 15287:2025, e NÃO a 14724.
%
%  Copie este arquivo, renomeie e escreva seu trabalho aqui mesmo: todas
%  as divisões do texto ficam neste único arquivo. Se preferir dividir o
%  texto em vários arquivos, veja o manual (manual.pdf) — ele próprio é
%  organizado assim e mostra como fazer.
%
%  Compile com:  latexmk modelo-tcc1.tex
% =====================================================================
\documentclass[12pt,a4paper,oneside]{book}

% O tipo governa o modelo inteiro: natureza, folha de aprovação, resumos e
% metadados exigidos. O TCC2 é outro documento, regido por outra norma —
% ao concluir o trabalho, comece de modelo-tcc2.tex em vez de trocar o
% valor abaixo.
\usepackage[tipo=tcc1]{UnifeiICTReport}

\usepackage{csquotes}
\usepackage{fontspec}

% --- Referências bibliográficas --------------------------------------
% Para citar trabalhos e produzir a lista de referências, crie o SEU
% arquivo .bib (por exemplo, referencias.bib, nesta mesma pasta) e
% descomente as duas linhas abaixo e o \printbibliography lá embaixo.
%
% ATENÇÃO: não aponte para o referencias.bib deste repositório. Ele é
% fictício — traz referências inventadas que servem apenas aos exemplos
% do manual. Um trabalho entregue com ele tem bibliografia falsa, e isso
% não salta aos olhos ao folhear o PDF.
%
%\usepackage[style=abnt]{biblatex}
%\bibliography{referencias.bib}

% --- Siglas e símbolos ------------------------------------------------
% Declare aqui as siglas do seu trabalho. Depois de declarar ao menos
% uma, descomente também o \printglossary correspondente lá embaixo.
% Sem nenhuma declaração, a lista sairia como uma página de título
% seguida de nada.
%
%\makeglossaries
%\newacronym{ict}{ICT}{Instituto de Ciências e Tecnologia}

%----------------------------------------------------------------------
% Dados da capa e da folha de rosto
%----------------------------------------------------------------------

\title{Título do seu projeto de pesquisa}

% Escreva só o texto do subtítulo, sem os dois-pontos: o modelo insere
% essa pontuação sozinho. Comente a linha se não houver subtítulo.
%\subtitle{Subtítulo, se houver}

\author{Seu Nome Completo}

\supervisor{Prof. Dr. Nome do Orientador}

%\cosupervisor{Prof. Dra. Nome da Coorientadora}

% Curso de graduação. Entra na natureza do trabalho.
\subject{Nome do seu Curso}

% --- Banca examinadora ------------------------------------------------
% No TCC1 a defesa é FACULTATIVA. Sem nenhum \bancamembro declarado, a
% folha de aprovação simplesmente não é emitida — nenhuma página em
% branco, nenhuma entrada no sumário.
%
% Se houver defesa, declare apenas os membros EXTERNOS: o orientador já
% entra pela linha \supervisor acima e não deve ser repetido aqui.
%\bancamembro{Prof. Dr. Nome do Membro}{Doutor em Área}{Unifei}
%\bancamembro{Prof. Dra. Nome da Membra}{Doutora em Área}{Unifei}

%\aprovacaodata{15 de dezembro de 2026}
%\notaaprovacao{9,5}

%----------------------------------------------------------------------
% Resumo
%----------------------------------------------------------------------

\abstract{%
    Escreva aqui o resumo do projeto: o problema, os objetivos, a
    metodologia prevista e os resultados esperados, em um único parágrafo.
}
\keywords{primeira palavra-chave; segunda; terceira.}

% Resumo em língua estrangeira. Opcional aqui — a NBR 15287:2025 não
% prevê resumo entre os elementos do projeto de pesquisa, nem em língua
% vernácula nem estrangeira. Descomente se quiser incluí-lo.
%\abstractseclang{%
%    Write the abstract in the foreign language here.
%}
%\keywordsseclang{first keyword; second; third.}

%======================================================================
\begin{document}

\frontmatter

\maketitle

% Só é emitida se houver \bancamembro declarado acima.
\folhaaprovacao

\makeabstracts

% --- Listas -----------------------------------------------------------
% Descomente cada lista DEPOIS de o trabalho ter o que listar. Chamada
% sem conteúdo produz uma página com o título e nada abaixo.
%\listoffigures
%\listoftables
%\listofquadros
%\listofgraficos
%\printglossary[type=\acronymtype,title={\acronymlistname}]
%\printglossary[type=symbols,title={\symbolslistname}]

\tableofcontents

\mainmatter

% Atenção ao vocabulário: você digita \chapter, e o documento imprime
% "Seção". As divisões do trabalho não se chamam capítulos, então o
% modelo rebaixa cada nível um degrau:
%   \chapter -> Seção          \subsection    -> Subsubseção
%   \section -> Subseção       \subsubsection -> Parágrafo
% O manual explica isso em detalhe.

% A NBR 15287:2025 §4.2.2 prescreve o que a parte textual do projeto deve
% conter. As divisões abaixo cobrem esses elementos; reorganize-as como
% preferir, desde que todos apareçam.

\chapter{Introdução}\label{cap:introducao}

Apresente o tema, contextualize-o e delimite o problema de pesquisa.

\section{Problema}\label{sec:problema}

Enuncie com precisão a pergunta que a pesquisa pretende responder.

\section{Hipóteses}\label{sec:hipoteses}

Formule as hipóteses, quando couberem ao tipo de pesquisa.

\section{Objetivos}\label{sec:objetivos}

Apresente o objetivo geral e, em seguida, os objetivos específicos.

\section{Justificativa}\label{sec:justificativa}

Explique por que o problema merece ser investigado.

\chapter{Referencial teórico}\label{cap:referencial}

Apresente os trabalhos e conceitos que fundamentam a pesquisa.

\chapter{Metodologia}\label{cap:metodologia}

Descreva como a pesquisa será conduzida.

\chapter{Recursos e cronograma}\label{cap:recursos}

Relacione os recursos necessários e distribua as etapas no tempo
disponível.

% Descomente junto com o bloco de bibliografia do preâmbulo.
%\printbibliography

\backmatter

% Apêndices são de sua autoria; anexos são documentos de terceiros.
% Descomente o grupo que for usar (veja o manual para os detalhes).
%\apendices
%\apendice{ape:rotulo}{Título do apêndice}
%
%\anexos
%\anexo{ane:rotulo}{Título do anexo}

\end{document}
